% =====================================================
% RAPPORT D'ANALYSE — CAMPAGNES MARKETING
% Template : sikatikenmogne/latex-report-template (MIT)
% =====================================================

\documentclass{internshipreport}

\reporttitle{Analyse des performances des campagnes marketing}
\reportsubtitle{De PostgreSQL à Python : exploration, analyse et visualisation de 200\,005 campagnes marketing}
\reporttype{RAPPORT D'ANALYSE DE DONNÉES}

\reportauthor{Manassé Lumu Bamanoa}
\studentid{}
\program{Data Analytics / Data Analysis}
\academicyear{2026}
\internshipperiod{Année 2021 (données) --- rédaction 2026}

\university{Data Analytics}
\department{Analyse de données}
\academicsupervisor{}
\supervisorposition{}

\company{PostgreSQL, SQL, Python, Pandas, Matplotlib, Streamlit}
\companydepartment{Pipeline analytique}
\companytutor{}
\tutorposition{}

\defensedate{2026}

% =====================================================
% PACKAGES ET CONFIGURATION DE BASE
% Les packages cœur sont déjà chargés par internshipreport.cls.
% Ici : ressources du projet uniquement.
% =====================================================

\usepackage{csquotes}
\addbibresource{content/backmatter/bibliography.bib}
\usetikzlibrary{arrows.meta,positioning,shapes.geometric}
\graphicspath{{assets/images/charts/}{assets/images/}{assets/logos/}}
\setlength{\headheight}{14.5pt}

% =====================================================
% SCHÉMA DE COULEURS PERSONNALISÉ POUR LE TUTORIEL
% =====================================================

% ===== COULEURS PRINCIPALES DU PROJET =====

% Couleurs de marque (personnalisables selon l'entreprise)
\definecolor{brandprimary}{RGB}{25,118,210}        % Bleu corporate moderne
\definecolor{brandsecondary}{RGB}{46,125,50}       % Vert professionnel
\definecolor{brandaccent}{RGB}{255,112,67}         % Orange accent dynamique

% Variantes des couleurs principales
\definecolor{brandprimarylight}{RGB}{144,202,249}  % Bleu clair
\definecolor{brandsecondarylight}{RGB}{129,199,132} % Vert clair
\definecolor{brandaccentlight}{RGB}{255,183,77}    % Orange clair

\definecolor{brandprimarydark}{RGB}{13,71,161}     % Bleu foncé
\definecolor{brandsecondarydark}{RGB}{27,94,32}    % Vert foncé
\definecolor{brandaccentdark}{RGB}{230,74,25}      % Orange foncé

% ===== COULEURS NEUTRES PROFESSIONNELLES =====

% Échelle de gris moderne
\definecolor{neutraldark}{RGB}{33,33,33}           % Noir corporate
\definecolor{neutralmedium}{RGB}{97,97,97}         % Gris moyen
\definecolor{neutrallight}{RGB}{158,158,158}       % Gris clair
\definecolor{neutrallighter}{RGB}{224,224,224}     % Gris très clair
\definecolor{neutralwhite}{RGB}{255,255,255}       % Blanc pur

% Couleurs de fond
\definecolor{backgroundprimary}{RGB}{250,250,250}  % Fond principal
\definecolor{backgroundsecondary}{RGB}{245,245,245} % Fond secondaire
\definecolor{backgroundcard}{RGB}{255,255,255}     % Fond cartes/boîtes

% ===== COULEURS SÉMANTIQUES =====

% États et notifications
\definecolor{successcolor}{RGB}{76,175,80}         % Succès (vert)
\definecolor{warningcolor}{RGB}{255,152,0}         % Avertissement (orange)
\definecolor{errorcolor}{RGB}{244,67,54}           % Erreur (rouge)
\definecolor{infocolor}{RGB}{33,150,243}           % Information (bleu)

% Variantes claires pour les fonds
\definecolor{successlight}{RGB}{232,245,233}       % Fond succès
\definecolor{warninglight}{RGB}{255,243,224}       % Fond avertissement
\definecolor{errorlight}{RGB}{255,235,238}         % Fond erreur
\definecolor{infolight}{RGB}{227,242,253}          % Fond information

% ===== COULEURS POUR LE CODE SOURCE =====

% Syntaxe highlighting moderne
\definecolor{codebackground}{RGB}{248,249,250}     % Fond code
\definecolor{codeborder}{RGB}{218,220,224}         % Bordure code
\definecolor{codekeyword}{RGB}{215,58,73}          % Mots-clés
\definecolor{codecomment}{RGB}{106,153,85}         % Commentaires
\definecolor{codestring}{RGB}{3,102,214}           % Chaînes
\definecolor{codenumber}{RGB}{0,92,197}            % Nombres
\definecolor{codefunction}{RGB}{121,93,163}        % Fonctions
\definecolor{codevariable}{RGB}{227,98,9}          % Variables

% ===== COULEURS POUR LES TABLEAUX =====

% Style moderne pour tableaux
\definecolor{tableheadercolor}{RGB}{37,99,235}     % En-têtes (bleu)
\definecolor{tableheadertext}{RGB}{255,255,255}    % Texte en-têtes
\definecolor{tablealternatecolor}{RGB}{248,250,252} % Lignes alternées
\definecolor{tablehovercolor}{RGB}{219,234,254}    % Survol (bleu clair)
\definecolor{tableborder}{RGB}{203,213,225}        % Bordures

% ===== COULEURS POUR LES GRAPHIQUES =====

% Palette accessible pour graphiques (respecte les standards d'accessibilité)
\definecolor{chartblue}{RGB}{54,162,235}           % Bleu graphique
\definecolor{chartred}{RGB}{255,99,132}            % Rouge graphique
\definecolor{chartgreen}{RGB}{75,192,192}          % Vert graphique
\definecolor{chartyellow}{RGB}{255,205,86}         % Jaune graphique
\definecolor{chartpurple}{RGB}{153,102,255}        % Violet graphique
\definecolor{chartorange}{RGB}{255,159,64}         % Orange graphique

% ===== COULEURS POUR LA GESTION DE PROJET =====

% États des tâches
\definecolor{todostatus}{RGB}{156,163,175}         % À faire (gris)
\definecolor{inprogressstatus}{RGB}{59,130,246}    % En cours (bleu)
\definecolor{reviewstatus}{RGB}{245,158,11}        % En révision (orange)
\definecolor{donestatus}{RGB}{34,197,94}           % Terminé (vert)
\definecolor{blockedstatus}{RGB}{239,68,68}        % Bloqué (rouge)

% Priorités
\definecolor{prioritylow}{RGB}{156,163,175}        % Basse (gris)
\definecolor{prioritymedium}{RGB}{245,158,11}      % Moyenne (orange)
\definecolor{priorityhigh}{RGB}{239,68,68}         % Haute (rouge)
\definecolor{prioritycritical}{RGB}{147,51,234}    % Critique (violet)

% ===== CONFIGURATION DES LIENS HYPERTEXTE =====

\hypersetup{
    colorlinks=true,
    linkcolor=neutraldark,          % Liens internes
    citecolor=brandsecondary,        % Citations
    filecolor=brandaccent,           % Fichiers
    urlcolor=brandprimary,           % URLs
    anchorcolor=brandprimarydark,    % Ancres
    menucolor=brandprimary,          % Menus PDF
    runcolor=brandaccent,            % Exécutables
    bookmarksnumbered=true,
    bookmarksopen=true,
    bookmarksopenlevel=1,
    pdfstartview=FitH,
    pdfpagelayout=OneColumn,
    pdfdisplaydoctitle=true
}

% ===== COULEURS POUR L'IMPRESSION =====

% Alternatives haute-contraste pour l'impression
\definecolor{printblack}{RGB}{0,0,0}               % Noir impression
\definecolor{printgray}{RGB}{64,64,64}             % Gris impression
\definecolor{printlightgray}{RGB}{128,128,128}     % Gris clair impression

% Commande pour basculer en mode impression
\newif\ifprintmode
\printmodefalse  % Mode écran par défaut

% Redéfinition conditionnelle des couleurs pour l'impression
\ifprintmode
    \definecolor{brandprimary}{RGB}{0,0,0}
    \definecolor{brandsecondary}{RGB}{64,64,64}
    \definecolor{brandaccent}{RGB}{128,128,128}
\fi

% =====================================================
% COMMANDES PERSONNALISÉES POUR LE TUTORIEL
% =====================================================

% ===== COMMANDES SPÉCIFIQUES AU PROJET =====

% Références d'entreprise et organisations
\newcommand{\companyref}[1]{\textcolor{primarycolor}{\textbf{#1}}}
\newcommand{\universityref}[1]{\textcolor{secondarycolor}{\textbf{#1}}}

% Nom du projet et technologies principales
\newcommand{\projectref}{\textbf{\color{brandprimary}\projectname}}
\newcommand{\mainframework}{\framework{\maintech}}
\newcommand{\backend}{\framework{\backendtech}}
\newcommand{\maindatabase}{\database{\databasetech}}
\newcommand{\programming}[1]{\texttt{\color{primarycolor}#1}}

% ===== COMMANDES POUR LES MÉTHODOLOGIES =====

% Méthodologies de développement
\newcommand{\agile}{\methodology{Agile}}
\newcommand{\scrum}{\methodology{Scrum}}
\newcommand{\kanban}{\methodology{Kanban}}
\newcommand{\devops}{\methodology{DevOps}}

% Éléments Scrum spécifiques
\newcommand{\sprintduration}[1]{\duration{#1-week sprint}}
\newcommand{\userstoryref}[1]{\textbf{\color{brandaccent}US-#1}}
\newcommand{\taskref}[1]{\textbf{\color{brandsecondary}TASK-#1}}

% ===== COMMANDES POUR LES TESTS =====

% Types de tests
\newcommand{\unittest}{\testcase{Unit Test}}
\newcommand{\integrationtest}{\testcase{Integration Test}}
% \newcommand{\e2etest}{\testcase{End-to-End Test}}
\newcommand{\performancetest}{\testcase{Performance Test}}

% Couverture de tests
\newcommand{\codecoverage}[1]{\metric{#1\% coverage}}
\newcommand{\testpass}[1]{\textcolor{brandsuccess}{\textbf{#1 passed}}}
\newcommand{\testfail}[1]{\textcolor{branddanger}{\textbf{#1 failed}}}

% ===== COMMANDES POUR LES PERFORMANCES =====

% Métriques de performance
\newcommand{\loadtime}[1]{\metric{#1ms load time}}
\newcommand{\responsetime}[1]{\metric{#1ms response}}
\newcommand{\throughput}[1]{\metric{#1 req/sec}}
\newcommand{\concurrent}[1]{\metric{#1 concurrent users}}

% Améliorations de performance
\newcommand{\speedup}[1]{\improvement{#1× faster}}
\newcommand{\reduction}[1]{\improvement{#1\% reduction}}
\newcommand{\increase}[1]{\improvement{#1\% increase}}

% ===== COMMANDES POUR LA SÉCURITÉ =====

% Éléments de sécurité
\newcommand{\authentication}{\important{Authentication}}
\newcommand{\authorization}{\important{Authorization}}
\newcommand{\encryption}{\important{Encryption}}
\newcommand{\validation}{\important{Input Validation}}

% Types d'utilisateurs et rôles
% \newcommand{\userrole}[1]{\role{#1}}
\newcommand{\adminrole}{\role{Administrator}}
\newcommand{\userrole}{\role{User}}
\newcommand{\guestrole}{\role{Guest}}

% ===== COMMANDES POUR LES APIs =====

% Méthodes HTTP
\newcommand{\httpget}{\httpmethod{GET}}
\newcommand{\httppost}{\httpmethod{POST}}
\newcommand{\httpput}{\httpmethod{PUT}}
\newcommand{\httpdelete}{\httpmethod{DELETE}}

% Codes de statut HTTP
\newcommand{\statusok}{\statuscode{200 OK}}
\newcommand{\statuscreated}{\statuscode{201 Created}}
\newcommand{\statusnotfound}{\statuscode{404 Not Found}}
\newcommand{\statuserror}{\statuscode{500 Internal Server Error}}

% Endpoints API
\newcommand{\apiendpoint}[1]{\endpoint{/api/v1/#1}}

% ===== COMMANDES POUR LA DOCUMENTATION =====

% Types de documentation
\newcommand{\apidoc}{\deliverable{API Documentation}}
\newcommand{\userdoc}{\deliverable{User Documentation}}
\newcommand{\techdoc}{\deliverable{Technical Documentation}}

% Outils de documentation
\newcommand{\swagger}{\software{Swagger}}
\newcommand{\postman}{\software{Postman}}
\newcommand{\jsdoc}{\software{JSDoc}}

% ===== COMMANDES POUR LES OUTILS =====

% Outils de développement
\newcommand{\vscode}{\software{Visual Studio Code}}
\newcommand{\git}{\software{Git}}
\newcommand{\github}{\software{GitHub}}
\newcommand{\docker}{\software{Docker}}
\newcommand{\kubernetes}{\software{Kubernetes}}

% Services cloud
\newcommand{\aws}{\software{Amazon Web Services}}
\newcommand{\azure}{\software{Microsoft Azure}}
\newcommand{\gcp}{\software{Google Cloud Platform}}

% ===== COMMANDES POUR LES DATES ET DURÉES =====

% Phases du projet
\newcommand{\phase}[1]{\milestone{Phase #1}}
\newcommand{\iteration}[1]{\milestone{Iteration #1}}
\newcommand{\release}[1]{\milestone{Release #1}}

% Durées typiques
\newcommand{\oneweek}{\duration{1 week}}
\newcommand{\twoweeks}{\duration{2 weeks}}
\newcommand{\onemonth}{\duration{1 month}}

% Références croisées (préfixes de labels du projet)
\newcommand{\filename}[1]{\texttt{#1}}
\newcommand{\figref}[1]{\figurename~\ref{fig:#1}}
\newcommand{\tabref}[1]{\tablename~\ref{tab:#1}}
\newcommand{\chapref}[1]{\chaptername~\ref{chap:#1}}
\newcommand{\secref}[1]{section~\ref{sec:#1}}

\newtcolorbox{codeexplanation}[1][]{
    enhanced,
    colback=codebackground,
    colframe=codeborder,
    boxrule=0.6pt,
    arc=2pt,
    left=10pt,
    right=10pt,
    top=8pt,
    bottom=8pt,
    fonttitle=\bfseries,
    title={Commentaire de code},
    #1
}

% ===== COMMANDES POUR LES RÉFÉRENCES =====

% Références aux livrables
\newcommand{\deliverableref}[1]{Deliverable~\ref{deliv:#1}}
\newcommand{\milestoneref}[1]{Milestone~\ref{mile:#1}}
\newcommand{\requirementref}[1]{Requirement~\ref{req:#1}}

% Références aux tests
\newcommand{\testcaseref}[1]{Test Case~\ref{test:#1}}
\newcommand{\testsuiteref}[1]{Test Suite~\ref{suite:#1}}

% ===== COMMANDES CONDITIONNELLES =====

% Compilation conditionnelle
\newif\ifshowdetails
\showdetailstrue  % Afficher les détails par défaut

\newcommand{\detailsonly}[1]{\ifshowdetails#1\fi}
\newcommand{\summaryonly}[1]{\ifshowdetails\else#1\fi}

% Version du document
\newcommand{\version}[1]{\textit{\color{neutralmedium}v#1}}
\newcommand{\updated}[1]{\textit{\color{neutralmedium}Updated: #1}}


% =====================================================
% CONFIGURATION AVANCÉE DU STYLE
% =====================================================

% Commandes \frontmatter / \mainmatter (classe book), absentes de report
% et utilisées par le template dans main.tex.
\makeatletter
\newif\if@mainmatter
\@mainmattertrue
\providecommand{\frontmatter}{%
  \cleardoublepage
  \@mainmatterfalse
  \pagenumbering{roman}%
}
\providecommand{\mainmatter}{%
  \cleardoublepage
  \@mainmattertrue
  \pagenumbering{arabic}%
}
\makeatother

% ===== TYPOGRAPHIE PERSONNALISÉE =====

% Interligne (choix: \singlespacing, \onehalfspacing, \doublespacing)
\onehalfspacing  % Interligne 1.5 (recommandé pour les rapports)

% Espacement des paragraphes
\setlength{\parskip}{8pt plus 2pt minus 1pt}    % 8pt entre paragraphes
\setlength{\parindent}{0pt}                     % Pas d'indentation première ligne

% Marges personnalisées (optionnel - déjà défini dans la classe)
% \geometry{
%     top=2.5cm,
%     bottom=2.5cm,
%     left=3cm,
%     right=2.5cm,
%     headheight=15pt
% }

% ===== PERSONNALISATION DES TITRES =====

% Espacement autour des titres
\titlespacing*{\chapter}{0pt}{-10pt}{25pt}      % Chapitre: espace avant/après
\titlespacing*{\section}{0pt}{18pt}{12pt}       % Section: espace avant/après
\titlespacing*{\subsection}{0pt}{14pt}{8pt}     % Sous-section: espace avant/après
\titlespacing*{\subsubsection}{0pt}{12pt}{6pt}  % Sous-sous-section

% Couleurs des titres (personnalisable)
\titleformat{\chapter}[display]
    {\normalfont\huge\bfseries\color{brandprimary}}
    {\chaptertitlename\ \thechapter}{20pt}{\Huge\MakeUppercase}

\titleformat{\section}
    {\normalfont\Large\bfseries\color{brandsecondary}}
    {\thesection}{1em}{}

\titleformat{\subsection}
    {\normalfont\large\bfseries\color{neutraldark}}
    {\thesubsection}{1em}{}

% ===== CONFIGURATION DES LISTES =====

% Listes à puces - niveau 1
\setlist[itemize,1]{
    label=\textcolor{neutraldark}{\textbullet},
    leftmargin=20pt,
    itemsep=3pt,
    parsep=1pt,
    topsep=6pt,
    partopsep=0pt
}

% Listes à puces - niveau 2
\setlist[itemize,2]{
    label=\textcolor{neutraldark}{\textendash},
    leftmargin=35pt,
    itemsep=2pt,
    parsep=0pt,
    topsep=3pt
}

% Listes numérotées
\setlist[enumerate,1]{
    leftmargin=20pt,
    itemsep=3pt,
    parsep=1pt,
    topsep=6pt
}

% ===== STYLE DES TABLEAUX =====

% Espacement global des tableaux
\renewcommand{\arraystretch}{1.4}  % Hauteur des lignes
\setlength{\tabcolsep}{10pt}       % Espacement des colonnes

% Style des en-têtes de tableaux
\renewcommand{\tableheadercell}[1]{%
    \cellcolor{tableheadercolor}\textbf{\color{neutraldark}#1}%
}

% ===== CONFIGURATION DES FIGURES =====

% Positionnement par défaut des figures
\renewcommand{\topfraction}{0.9}       % Max 90% de la page pour figures en haut
\renewcommand{\bottomfraction}{0.9}    % Max 90% de la page pour figures en bas
\renewcommand{\textfraction}{0.1}      % Min 10% de la page pour le texte
\renewcommand{\floatpagefraction}{0.8} % Min 80% pour page de figures

% Style des légendes
\captionsetup[figure]{
    font={small},
    labelfont={bf,color=neutraldark},
    textfont={color=neutraldark},
    margin=20pt,
    skip=12pt,
    position=bottom
}

\captionsetup[table]{
    font={small},
    labelfont={bf,color=neutraldark},
    textfont={color=neutraldark},
    margin=20pt,
    skip=8pt,
    position=top
}

% ===== STYLE DU CODE SOURCE =====

% Configuration globale pour tous les listings
\lstset{
    backgroundcolor=\color{codebackground},
    basicstyle=\footnotesize\ttfamily\color{neutraldark},
    commentstyle=\color{codecomment}\itshape,
    keywordstyle=\color{codekeyword}\bfseries,
    stringstyle=\color{codestring},
    numberstyle=\tiny\color{neutralmedium},
    numbers=left,
    numbersep=12pt,
    frame=single,
    rulecolor=\color{codeborder},
    frameround=tttt,
    framesep=10pt,
    xleftmargin=20pt,
    xrightmargin=10pt,
    breaklines=true,
    breakatwhitespace=false,
    tabsize=2,
    showspaces=false,
    showstringspaces=false,
    showtabs=false,
    captionpos=b
}

% ===== EN-TÊTES ET PIEDS DE PAGE =====

% Style pour les pages normales
\fancypagestyle{mainmatter}{
    \fancyhf{}
    \fancyhead[C]{\footnotesize\textcolor{neutralmedium}{\MakeUppercase{\@reporttitle}}}
    \fancyfoot[R]{\footnotesize\textcolor{neutralmedium}{%
        \leftmark \quad \textbf{|} \quad \textbf{\thepage}%
    }}
    \renewcommand{\headrulewidth}{0pt}
    \renewcommand{\footrulewidth}{0.5pt}
    \renewcommand{\footrule}{\color{neutrallight}\hrule width\headwidth height\footrulewidth}
}

% Style pour les pages de chapitre
\fancypagestyle{plain}{
    \fancyhf{}
    \fancyfoot[C]{\textcolor{neutraldark}{\thepage}}
    \renewcommand{\headrulewidth}{0pt}
    \renewcommand{\footrulewidth}{0pt}
}

% ===== PERSONNALISATION TABLE DES MATIÈRES =====

% Style du titre de la table des matières
% \renewcommand{\cfttoctitlefont}{\Large\bfseries\color{brandprimary}}

% Style des entrées de chapitre
\renewcommand{\cftchapfont}{\bfseries\color{neutraldark}}
\renewcommand{\cftchappagefont}{\bfseries\color{neutraldark}}
\renewcommand{\cftchapleader}{\cftdotfill{\cftsecdotsep}}

% Style des entrées de section
\renewcommand{\cftsecfont}{\color{neutraldark}}
\renewcommand{\cftsecpagefont}{\color{neutralmedium}}

% =====================================================
% MÉTADONNÉES DU RAPPORT D'ANALYSE
% =====================================================

\renewcommand{\reporttitle}{Analyse des performances des campagnes marketing}
\renewcommand{\reportsubtitle}{De PostgreSQL à Python : exploration, analyse et visualisation de 200\,005 campagnes marketing}
\renewcommand{\reporttype}{RAPPORT D'ANALYSE DE DONNÉES}

\renewcommand{\reportauthor}{Manassé Lumu Bamanoa}
\renewcommand{\studentid}{}
\renewcommand{\program}{Data Analytics / Data Analysis}
\renewcommand{\academicyear}{2026}
\renewcommand{\internshipperiod}{Année 2021 (données) --- rédaction 2026}

\renewcommand{\university}{Data Analytics}
\renewcommand{\department}{Analyse de données}
\renewcommand{\academicsupervisor}{}
\renewcommand{\supervisorposition}{}

\renewcommand{\company}{PostgreSQL, SQL, Python, Pandas, Matplotlib, Streamlit}
\renewcommand{\companydepartment}{Pipeline analytique}
\renewcommand{\companytutor}{}
\renewcommand{\tutorposition}{}

\renewcommand{\defensedate}{2026}

\newcommand{\projectname}{Marketing Campaigns Analysis}
\newcommand{\maintech}{Python}
\newcommand{\backendtech}{PostgreSQL}
\newcommand{\databasetech}{PostgreSQL}

\hypersetup{
    pdftitle={Analyse des performances des campagnes marketing},
    pdfauthor={Manassé Lumu Bamanoa},
    pdfsubject={Data Analytics / Data Analysis},
    pdfkeywords={Data Analytics, PostgreSQL, SQL, Python, Pandas, Streamlit}
}


% =====================================================
% ENVIRONNEMENTS PERSONNALISÉS
% =====================================================

% Boîte d'information
\newtcolorbox{infobox}[1][Information]{
    enhanced,
    colback=primarycolor!10,
    colframe=primarycolor,
    boxrule=1pt,
    arc=3pt,
    left=10pt,
    right=10pt,
    top=8pt,
    bottom=8pt,
    fonttitle=\bfseries\color{white},
    title={#1}
}

% Boîte d'avertissement
\newtcolorbox{warningbox}[1][Attention]{
    enhanced,
    colback=orange!10,
    colframe=orange,
    boxrule=1pt,
    arc=3pt,
    left=10pt,
    right=10pt,
    top=8pt,
    bottom=8pt,
    fonttitle=\bfseries\color{white},
    title={#1}
}

% Boîte de succès
\newtcolorbox{successbox}[1][Résultat]{
    enhanced,
    colback=secondarycolor!10,
    colframe=secondarycolor,
    boxrule=1pt,
    arc=3pt,
    left=10pt,
    right=10pt,
    top=8pt,
    bottom=8pt,
    fonttitle=\bfseries\color{white},
    title={#1}
}

% Environnement pour objectifs
% \newenvironment{objectives}
% {
%     \par\vspace{10pt}
%     \noindent\textbf{\color{primarycolor}Objectives:}
%     \par\vspace{5pt}
% }
% {
%     \par\vspace{5pt}
% }

% Environnement pour résultats
% \newenvironment{results}
% {
%     \par\vspace{10pt}
%     \noindent\textbf{\color{secondarycolor}Results:}
%     \par\vspace{5pt}
% }
% {
%     \par\vspace{5pt}
% }

% Environnement pour définitions
% \newenvironment{definition}[1]
% {
%     \par\vspace{10pt}
%     \noindent\textbf{\color{primarycolor}Definition - #1:}
%     \par\vspace{5pt}
%     \begin{quote}
% }
% {
%     \end{quote}
%     \par\vspace{5pt}
% }

% Missing environments
\newtcolorbox{technicalbox}[1][Note technique]{
    enhanced,
    colback=blue!5,
    colframe=blue!50!black,
    boxrule=1pt,
    arc=3pt,
    left=10pt,
    right=10pt,
    top=8pt,
    bottom=8pt,
    fonttitle=\bfseries\color{white},
    title={#1}
}

\newtcolorbox{bestpractice}[1][Bonne pratique]{
    enhanced,
    colback=green!5,
    colframe=green!50!black,
    boxrule=1pt,
    arc=3pt,
    left=10pt,
    right=10pt,
    top=8pt,
    bottom=8pt,
    fonttitle=\bfseries\color{white},
    title={#1}
}

\newtcolorbox{lessonslearned}[1][Retour d'expérience]{
    enhanced,
    colback=purple!5,
    colframe=purple!50!black,
    boxrule=1pt,
    arc=3pt,
    left=10pt,
    right=10pt,
    top=8pt,
    bottom=8pt,
    fonttitle=\bfseries\color{white},
    title={#1}
}

% =====================================================
% TEMPLATES DE FIGURES
% =====================================================

% Template de figure standard (légende en bas)
\newcommand{\standardfigure}[4]{
  \begin{figure}[H]
    \centering
    \includegraphics[width=#2, keepaspectratio]{#1}
    \caption{#3}
    \label{#4}
  \end{figure}
}

% Template de figure avec cadre
\newcommand{\framedfigure}[4]{
  \begin{figure}[H]
    \centering
    \fbox{\includegraphics[width=#2, keepaspectratio]{#1}}
    \caption{#3}
    \label{#4}
  \end{figure}
}

% Template de figure simple sans cadre
\newcommand{\simplefigure}[4]{
  \begin{figure}[h]
    \centering
    \includegraphics[width=#2]{#1}
    \caption{#3}
    \label{#4}
  \end{figure}
}

% Template pour diagrammes avec bordure colorée
\newcommand{\diagramfigure}[5]{
  \begin{figure}[H]
    \centering
    \begin{tcolorbox}[colback=white,colframe=primarycolor,width=#2]
      \centering
      \includegraphics[width=0.9\textwidth, height=#3, keepaspectratio]{#1}
    \end{tcolorbox}
    \caption{#4}
    \label{#5}
  \end{figure}
}

% Template pour figures côte à côte
\newcommand{\doublefigure}[8]{
  \begin{figure}[H]
    \centering
    \begin{subfigure}[b]{#3}
      \centering
      \includegraphics[width=\textwidth]{#1}
      \caption{#4}
      \label{#5}
    \end{subfigure}
    \hfill
    \begin{subfigure}[b]{#6}
      \centering
      \includegraphics[width=\textwidth]{#2}
      \caption{#7}
      \label{#8}
    \end{subfigure}
    \caption{Comparison figures}
  \end{figure}
}

% =====================================================
% TEMPLATES DE TABLEAUX
% =====================================================

% Template de tableau standard
\newenvironment{standardtable}[2]
{
  \begin{table}[h]
  \centering
  \rowcolors{2}{tablealternate}{white}
  \begin{tabular}{#1}
  \rowcolor{tableheader}
  #2 \\
  \hline
}
{
  \hline
  \end{tabular}
  \end{table}
}

% Template de tableau avec alternance
\newenvironment{alternatingtable}[2]
{
  \begin{table}[h]
  \centering
  \rowcolors{2}{tablealternate}{white}
  \begin{tabular}{#1}
  \rowcolor{tableheader}
  #2 \\
  \hline
}
{
  \hline
  \end{tabular}
  \end{table}
}

% Template de tableau coloré
\newenvironment{coloredtable}[3]
{
  \begin{table}[h]
  \centering
  \rowcolors{2}{#3}{white}
  \begin{tabular}{#1}
  \rowcolor{tableheader}
  #2 \\
  \hline
}
{
  \hline
  \end{tabular}
  \end{table}
}

% Template pour tableau de métriques
\newcommand{\metricstable}[1]{
  \begin{table}[h]
  \centering
  \begin{tabular}{|>{\columncolor{tableheader}}p{5cm}|>{\centering\arraybackslash}p{2cm}|>{\centering\arraybackslash}p{2cm}|>{\centering\arraybackslash\color{textcolor}\bfseries}p{2cm}|}
  \hline
  \rowcolor{tableheader}
  \textbf{Indicator} & \textbf{Before} & \textbf{After} & \textbf{Improvement} \\
  \hline
  #1
  \end{tabular}
  \end{table}
}

% Template de tableau simple sans bordures
\newenvironment{simpletable}[2]
{
  \begin{table}[h]
  \centering
  \begin{tabular}{#1}
  \toprule
  #2 \\
  \midrule
}
{
  \bottomrule
  \end{tabular}
  \end{table}
}


\begin{document}

% =====================================================
% PAGE DE GARDE — ÉTUDE PROFESSIONNELLE
% (bordure TikZ du template conservée ; contenu non académique)
% =====================================================

\begin{titlepage}
    \newgeometry{top=1.27cm,bottom=1.27cm,left=1.27cm,right=1.27cm}
    \thispagestyle{empty}

    {\singlespacing

    \begin{tikzpicture}[remember picture,overlay]
        \def\bordermarg{1cm}
        \def\innerspace{0.3cm}

        \draw[primarycolor, line width=2pt]
            ([xshift=\bordermarg,yshift=-\bordermarg]current page.north west)
            rectangle
            ([xshift=-\bordermarg,yshift=\bordermarg]current page.south east);

        \draw[primarycolor, line width=0.8pt]
            ([xshift=\bordermarg+\innerspace,yshift=-\bordermarg-\innerspace]current page.north west)
            rectangle
            ([xshift=-\bordermarg-\innerspace,yshift=\bordermarg+\innerspace]current page.south east);

        \draw[primarycolor, line width=1.5pt]
            ([xshift=\bordermarg,yshift=-\bordermarg]current page.north west)
            -- ++(1.2,0) -- ++(0,-0.4) -- ++(-0.4,0) -- ++(0,-1.2) -- ++(-0.8,0) -- cycle;
        \draw[primarycolor, line width=1pt, rotate around={45:([xshift=\bordermarg+0.6cm,yshift=-\bordermarg-0.6cm]current page.north west)}]
            ([xshift=\bordermarg+0.6cm,yshift=-\bordermarg-0.6cm]current page.north west) rectangle ++(0.15,0.15);

        \draw[primarycolor, line width=1.5pt]
            ([xshift=-\bordermarg,yshift=-\bordermarg]current page.north east)
            -- ++(-1.2,0) -- ++(0,-0.4) -- ++(0.4,0) -- ++(0,-1.2) -- ++(0.8,0) -- cycle;
        \draw[primarycolor, line width=1pt, rotate around={45:([xshift=-\bordermarg-0.6cm,yshift=-\bordermarg-0.6cm]current page.north east)}]
            ([xshift=-\bordermarg-0.6cm,yshift=-\bordermarg-0.6cm]current page.north east) rectangle ++(0.15,0.15);

        \draw[primarycolor, line width=1.5pt]
            ([xshift=\bordermarg,yshift=\bordermarg]current page.south west)
            -- ++(1.2,0) -- ++(0,0.4) -- ++(-0.4,0) -- ++(0,1.2) -- ++(-0.8,0) -- cycle;
        \draw[primarycolor, line width=1pt, rotate around={45:([xshift=\bordermarg+0.6cm,yshift=\bordermarg+0.6cm]current page.south west)}]
            ([xshift=\bordermarg+0.6cm,yshift=\bordermarg+0.6cm]current page.south west) rectangle ++(0.15,0.15);

        \draw[primarycolor, line width=1.5pt]
            ([xshift=-\bordermarg,yshift=\bordermarg]current page.south east)
            -- ++(-1.2,0) -- ++(0,0.4) -- ++(0.4,0) -- ++(0,1.2) -- ++(0.8,0) -- cycle;
        \draw[primarycolor, line width=1pt, rotate around={45:([xshift=-\bordermarg-0.6cm,yshift=\bordermarg+0.6cm]current page.south east)}]
            ([xshift=-\bordermarg-0.6cm,yshift=\bordermarg+0.6cm]current page.south east) rectangle ++(0.15,0.15);
    \end{tikzpicture}

    \begin{center}
        \vspace*{2.1cm}

        {\large \textcolor{textcolor}{\textbf{Data Analytics / Data Analysis}}}\\[1.35cm]

        {\Huge \textcolor{primarycolor}{\textbf{ANALYSE DES PERFORMANCES}}}\\[0.28cm]
        {\Huge \textcolor{primarycolor}{\textbf{DES CAMPAGNES MARKETING}}}\\[0.75cm]

        \textcolor{primarycolor}{\rule{0.62\linewidth}{0.5mm}}\\[0.7cm]

        {\large \textcolor{secondarycolor}{Exploration, analyse et visualisation}}\\[0.12cm]
        {\large \textcolor{secondarycolor}{de 200\,005 campagnes marketing}}\\[1.6cm]

        {\large \textcolor{textcolor}{Réalisé par}}\\[0.4cm]
        {\LARGE \textcolor{primarycolor}{\textbf{Manassé Lumu Bamanoa}}}\\[0.45cm]
        {\large \textcolor{graycolor}{Data Analytics / Data Analysis}}

        \vfill

        {\normalsize \textcolor{textcolor}{Étude réalisée le 21 septembre 2026}}\\[0.85cm]

        {\large \textcolor{textcolor}{\textbf{MediJr}}}\\[0.22cm]
        {\large \href{https://github.com/MediJr}{\textcolor{primarycolor}{\texttt{https://github.com/MediJr}}}}

        \vspace*{1.55cm}
    \end{center}

    } % singlespacing
    \restoregeometry
\end{titlepage}


\newgeometry{top=1.75cm,bottom=1.75cm,left=2.54cm,right=2.54cm}
\frontmatter

\chapter*{Résumé}
\addcontentsline{toc}{chapter}{Résumé}

Ce rapport documente un projet d'analyse de performances de campagnes marketing, réalisé sur un jeu de \metric{200\,005} campagnes observées du 1\textsuperscript{er} janvier au 31 décembre 2021. L'objectif est d'explorer, de décrire et de restituer les indicateurs de performance (ROI, coût d'acquisition, conversion, CTR, engagement) selon les entreprises, types de campagnes, canaux, audiences, segments, localisations et la période, sans inférence causale.

Les données sources proviennent d'un fichier Excel, chargées dans PostgreSQL (\database{marketing\_analysis}, table \filename{public.marketing\_campaigns}), explorées en SQL, extraites vers \technology{Pandas}, profilées, analysées (EDA), visualisées avec \technology{Matplotlib} puis restituées dans un dashboard \technology{Streamlit}. Le stockage repose sur une table relationnelle unique, volontairement sans modèle en étoile ni en flocon.

Les résultats principaux sont descriptifs : volumes très équilibrés entre modalités ; rythme d'environ 548 campagnes par jour civil ; ROI moyen de 5,0024, stable d'une dimension à l'autre ; CTR moyen de 14,04~\% contre une médiane de 9,98~\% ; seules corrélations notables entre CTR, clics et impressions, liées à la définition du ratio. Aucune anomalie de qualité n'a été détectée.

\textbf{Mots-clés :} Data Analytics, PostgreSQL, SQL, Python, Pandas, Data Visualization, Streamlit, Marketing Analytics.

\newpage
\chapter*{Table des matières}
\vspace*{-2cm}
\addcontentsline{toc}{chapter}{Table des matières}
\begin{singlespace}
\renewcommand{\contentsname}{}
\vspace*{-1cm}
\tableofcontents
\end{singlespace}

\newpage
\chapter*{Liste des figures}
\vspace*{-2cm}
\addcontentsline{toc}{chapter}{Liste des figures}
\begin{singlespace}
\renewcommand{\listfigurename}{}
\vspace*{-1cm}
\listoffigures
\end{singlespace}

\newpage
\chapter*{Liste des tableaux}
\vspace*{-2cm}
\addcontentsline{toc}{chapter}{Liste des tableaux}
\begin{singlespace}
\renewcommand{\listtablename}{}
\vspace*{-1cm}
\listoftables
\end{singlespace}

\chapter*{Liste des abréviations}
\addcontentsline{toc}{chapter}{Liste des abréviations}

\begin{table}[H]
\centering
\begin{tabular}{|p{3.2cm}|p{10.5cm}|}
\hline
\tableheadercell{Abréviation} & \tableheadercell{Définition} \\
\hline
CTR & Click-Through Rate (clics / impressions) \\
\hline
\alternaterow
EDA & Exploratory Data Analysis (analyse exploratoire) \\
\hline
KPI & Key Performance Indicator (indicateur de performance) \\
\hline
\alternaterow
ROI & Return on Investment (retour sur investissement) \\
\hline
SQL & Structured Query Language \\
\hline
\alternaterow
SGBD & Système de gestion de bases de données \\
\hline
\end{tabular}
\caption{Abréviations utilisées dans le rapport}
\label{tab:abbreviations}
\end{table}

\mainmatter

\chapter{Introduction}
\label{chap:introduction}

\section{Contexte}
\label{sec:contexte}

Le projet porte sur l'analyse de performances d'un ensemble de campagnes marketing. Les données sont disponibles sous forme d'un fichier Excel (\filename{marketing\_campaign\_dataset.xlsx}), puis stockées dans PostgreSQL et analysées en SQL et en Python. Le travail documenté ici reprend exclusivement les étapes réellement réalisées : import, exploration SQL, extraction Pandas, profilage, analyse exploratoire, visualisations Matplotlib, revue des graphiques et dashboard Streamlit \cite{pandasdocs,matplotlibdocs,streamlitdocs,postgresql18}.

Aucune hypothèse de marché, de budget publicitaire ou d'organisation interne n'est ajoutée : le rapport s'en tient aux variables et aux résultats observés dans le jeu de données \cite{marketingdataset}.

\section{Problématique}

Comment décrire, de façon reproductible et honnête, les performances de 200\,005 campagnes marketing de l'année 2021, selon les entreprises, les types de campagnes, les canaux, les audiences, les segments, les localisations et la période, sans transformer une association statistique en relation causale ?

\section{Objectif général}

Construire une chaîne d'analyse allant du fichier Excel jusqu'à une restitution interactive, en conservant l'intégrité des données sources et en produisant des indicateurs et des graphiques validés.

\section{Objectifs spécifiques}

\begin{objectives}
\begin{itemize}
\item Stocker les 200\,005 campagnes dans PostgreSQL, dans une table unique contrainte.
\item Explorer les volumes et les KPI en SQL (agrégations, \texttt{GROUP BY}, \texttt{HAVING}, \texttt{CASE}).
\item Extraire le résultat vers Pandas, contrôler les types et préparer des indicateurs calendaires et le CTR.
\item Réaliser une EDA exhaustive sans graphique, puis 16 visualisations Matplotlib revues.
\item Restituer les résultats dans un dashboard Streamlit filtrable, avec export CSV du sous-ensemble affiché.
\end{itemize}
\end{objectives}

\section{Périmètre}

Le périmètre est le fichier Excel du projet, la base \database{marketing\_analysis}, la table \filename{public.marketing\_campaigns} et les scripts Python associés. La période couverte par les dates de campagne est le 1\textsuperscript{er} janvier 2021 au 31 décembre 2021. Le grain est une ligne = une campagne.

Hors périmètre : Power BI, Tableau, Seaborn, Plotly, modélisation prédictive, modèle dimensionnel en étoile ou en flocon, et toute modification des données sources après import.

\section{Organisation du rapport}

Le \chapref{dataset} décrit le jeu de données. Le \chapref{architecture} présente le pipeline. Les \chapref{sql} et \chapref{python} documentent le stockage SQL puis l'extraction Python. Le \chapref{eda} synthétise l'EDA, le \chapref{visualisations} les graphiques validés et le \chapref{dashboard} le dashboard. Les \chapref{resultats} et \chapref{limites} rassemblent constats et limites, avant la conclusion et les annexes.

\chapter{Présentation du jeu de données}
\label{chap:dataset}

\section{Vue d'ensemble}

Le fichier source Excel contient \metric{200\,005} campagnes et \metric{15} colonnes. Après import PostgreSQL, la table conserve ces 15 colonnes. Après préparation Pandas, le DataFrame d'analyse en compte 21 (ajout de \texttt{year}, \texttt{month}, \texttt{month\_name}, \texttt{quarter}, \texttt{ctr} et \texttt{ctr\_percent}).

\begin{table}[H]
\centering
\caption{Caractéristiques générales du jeu de données}
\label{tab:dataset-overview}
\begin{tabular}{|p{6.2cm}|p{7.5cm}|}
\hline
\tableheadercell{Caractéristique} & \tableheadercell{Valeur observée} \\
\hline
Nombre de campagnes & 200\,005 \\
\hline
\alternaterow
Colonnes originales & 15 \\
\hline
Période des dates de campagne & 2021-01-01 $\rightarrow$ 2021-12-31 \\
\hline
\alternaterow
Entreprises & 5 \\
\hline
Types de campagnes & 5 \\
\hline
\alternaterow
Canaux & 6 \\
\hline
Audiences cibles & 5 \\
\hline
\alternaterow
Segments clients & 5 \\
\hline
Localisations & 5 \\
\hline
\end{tabular}
\end{table}

\section{Modalités catégorielles}

Les modalités suivantes ont été observées dans PostgreSQL et confirmées par l'EDA.

\begin{itemize}
\item \textbf{Entreprises :} TechCorp, Alpha Innovations, DataTech Solutions, NexGen Systems, Innovate Industries.
\item \textbf{Types :} Influencer, Search, Display, Email, Social Media.
\item \textbf{Canaux :} Email, Google Ads, YouTube, Instagram, Website, Facebook.
\item \textbf{Audiences :} Men 18-24, Men 25-34, Women 25-34, Women 35-44, All Ages.
\item \textbf{Segments :} Fashionistas, Foodies, Health \& Wellness, Outdoor Adventurers, Tech Enthusiasts.
\item \textbf{Localisations :} Miami, New York, Chicago, Los Angeles, Houston.
\end{itemize}

Les volumes sont quasi uniformes (environ 20~\% pour les variables à cinq modalités, environ 16,7~\% pour les six canaux). Toutes les combinaisons entreprise $\times$ canal (30), type $\times$ canal (30) et audience $\times$ segment (25) sont présentes.

\section{Variables principales}

Les indicateurs numériques utilisés tout au long du projet sont le ROI, le coût d'acquisition, le taux de conversion, les clics, les impressions, le score d'engagement, la durée et la date de campagne. Le CTR n'est pas une colonne source : il est calculé après extraction, campagne par campagne, comme le rapport clics / impressions.

\section{Dictionnaire de données}

\begin{table}[H]
\centering
\caption{Dictionnaire de données synthétique (colonnes sources)}
\label{tab:data-dictionary}
\scriptsize
\begin{tabular}{|p{3.1cm}|p{2.4cm}|p{7.6cm}|}
\hline
\tableheadercell{Colonne} & \tableheadercell{Type PG} & \tableheadercell{Rôle} \\
\hline
\texttt{campaign\_id} & \texttt{INTEGER} PK & Identifiant unique de campagne (1 à 200\,005) \\
\hline
\alternaterow
\texttt{company} & \texttt{TEXT} & Entreprise à l'origine de la campagne \\
\hline
\texttt{campaign\_type} & \texttt{TEXT} & Type de campagne \\
\hline
\alternaterow
\texttt{target\_audience} & \texttt{TEXT} & Audience ciblée \\
\hline
\texttt{duration\_days} & \texttt{INTEGER} & Durée en jours (15, 30, 45 ou 60) \\
\hline
\alternaterow
\texttt{channel\_used} & \texttt{TEXT} & Canal utilisé \\
\hline
\texttt{conversion\_rate} & \texttt{NUMERIC(4,2)} & Taux de conversion (0,01 à 0,15) \\
\hline
\alternaterow
\texttt{acquisition\_cost} & \texttt{NUMERIC(12,2)} & Coût d'acquisition de la campagne \\
\hline
\texttt{roi} & \texttt{NUMERIC(4,2)} & Retour sur investissement \\
\hline
\alternaterow
\texttt{location} & \texttt{TEXT} & Localisation \\
\hline
\texttt{campaign\_date} & \texttt{DATE} & Date de campagne \\
\hline
\alternaterow
\texttt{clicks} & \texttt{INTEGER} & Nombre de clics \\
\hline
\texttt{impressions} & \texttt{INTEGER} & Nombre d'impressions \\
\hline
\alternaterow
\texttt{engagement\_score} & \texttt{SMALLINT} & Score d'engagement (1 à 10) \\
\hline
\texttt{customer\_segment} & \texttt{TEXT} & Segment client \\
\hline
\end{tabular}
\end{table}

La colonne Excel \texttt{Duration} (texte du type \texttt{"15 days"}) a été convertie en entier de jours à l'import. La colonne \texttt{Date} a été convertie en type \texttt{DATE}.

\chapter{Architecture et pipeline}
\label{chap:architecture}

\section{Chaîne technique}

Le projet suit une chaîne linéaire, chaque étape consommant le résultat de la précédente sans réécrire PostgreSQL ni le fichier Excel une fois l'import effectué.

\begin{figure}[H]
\centering
\begin{tikzpicture}[
    node distance=7mm,
    box/.style={
        draw=primarycolor,
        line width=0.7pt,
        rounded corners=1pt,
        align=center,
        font=\small,
        inner sep=5pt,
        minimum width=7.4cm,
        fill=primarycolor!6
    },
    arr/.style={-{Stealth}, thick, primarycolor}
]
\node[box] (excel) {Excel\\ \texttt{marketing\_campaign\_dataset.xlsx}};
\node[box, below=of excel] (pg) {PostgreSQL\\ base \texttt{marketing\_analysis}, table unique};
\node[box, below=of pg] (sql) {SQL\\ exploration, agrégations, KPI};
\node[box, below=of sql] (py) {Python / Pandas\\ extraction lecture seule};
\node[box, below=of py] (prof) {Profiling / préparation\\ date, trimestre, CTR};
\node[box, below=of prof] (eda) {EDA\\ \texttt{python\_eda.py}};
\node[box, below=of eda] (mpl) {Matplotlib\\ 16 visualisations validées};
\node[box, below=of mpl] (st) {Streamlit\\ \texttt{dashboard.py}};
\draw[arr] (excel) -- (pg);
\draw[arr] (pg) -- (sql);
\draw[arr] (sql) -- (py);
\draw[arr] (py) -- (prof);
\draw[arr] (prof) -- (eda);
\draw[arr] (eda) -- (mpl);
\draw[arr] (mpl) -- (st);
\end{tikzpicture}
\caption{Pipeline analytique du projet (table relationnelle unique)}
\label{fig:pipeline}
\end{figure}

\section{Rôle de chaque étape}

\begin{enumerate}
\item \textbf{Excel.} Jeu source, non modifié après import.
\item \textbf{PostgreSQL.} Stockage structuré, contraintes, point de vérité pour les analyses SQL \cite{postgresql18}.
\item \textbf{SQL.} Lectures seules : volumes, moyennes, \texttt{CASE}, \texttt{HAVING}, analyses croisées et temporelles.
\item \textbf{Python / Pandas.} Extraction \texttt{SELECT *} ordonnée par \texttt{campaign\_id}, inspection des types et de la mémoire \cite{pandasdocs,sqlalchemydocs}.
\item \textbf{Profiling.} Contrôle des nulls, unicité de \texttt{campaign\_id}, conversion de la date, création des colonnes calendaires et du CTR.
\item \textbf{EDA.} Analyses A à S dans Pandas, sans graphique.
\item \textbf{Matplotlib.} Seize figures PNG, puis revue : seule l'échelle du rythme quotidien du graphique mensuel a été corrigée \cite{matplotlibdocs}.
\item \textbf{Streamlit.} Restitution filtrable des KPI et graphiques, export CSV du sous-ensemble affiché \cite{streamlitdocs}.
\end{enumerate}

\begin{infobox}[Choix d'architecture]
Le projet utilise volontairement \important{une table relationnelle unique}. Il n'y a pas de modèle en étoile, ni en flocon, ni de tables de faits et de dimensions séparées.
\end{infobox}

\chapter{Stockage et exploration SQL}
\label{chap:sql}

\section{Base et table}

La base \database{marketing\_analysis} a été créée sur PostgreSQL local (hôte \texttt{127.0.0.1}, port 5432). Le rôle d'application est \texttt{data\_project}. La table \filename{public.marketing\_campaigns} reprend les 15 colonnes source, en \texttt{snake\_case}.

Les types retenus sont : \texttt{INTEGER} pour l'identifiant, les clics, les impressions et la durée ; \texttt{SMALLINT} pour l'engagement ; \texttt{NUMERIC} pour le ROI, le coût et la conversion ; \texttt{DATE} pour la date ; \texttt{TEXT} pour les catégories. Toutes les colonnes sont \texttt{NOT NULL}. \texttt{campaign\_id} est clé primaire.

Contraintes \texttt{CHECK} posées à la création :

\begin{itemize}
\item \texttt{acquisition\_cost >= 0} ;
\item \texttt{conversion\_rate BETWEEN 0 AND 1} ;
\item \texttt{clicks <= impressions} ;
\item \texttt{engagement\_score BETWEEN 1 AND 10}.
\end{itemize}

\section{Import et contrôles}

Les 200\,005 lignes ont été importées depuis Excel (conversion de la durée texte vers un entier de jours, et de la date vers \texttt{DATE}). Les contrôles de qualité en base et en Pandas donnent : 0 valeur manquante, identifiants uniques de 1 à 200\,005, dates strictement dans l'année 2021, aucune anomalie bloquante.

\section{Analyses SQL}

Les requêtes d'exploration sont des \texttt{SELECT} en lecture seule. Elles ont servi à :

\begin{itemize}
\item \textbf{Agrégations et \texttt{GROUP BY}.} Volumes et ROI moyen par type, entreprise, canal, audience, segment, localisation, mois.
\item \textbf{KPI.} \texttt{COUNT}, \texttt{AVG} sur ROI, coût, conversion, engagement ; CTR calculé soit comme moyenne des ratios individuels, soit comme $\sum$ clics / $\sum$ impressions (ces deux définitions ne coïncident pas, d'où l'écart entre CTR moyen $\approx$ 14~\% et CTR global $\approx$ 10~\%).
\item \textbf{\texttt{CASE}.} Classes de ROI (faible $<$ 4, moyen $[4,6[$, élevé $\geq$ 6) : 66\,503 / 66\,438 / 67\,064 campagnes, ROI moyens 2,99 / 5,00 / 7,00.
\item \textbf{\texttt{HAVING}.} Filtres après agrégation, par exemple pour n'afficher que des groupes dépassant un seuil d'effectif ou de KPI.
\item \textbf{Analyses croisées.} Tous les couples entreprise $\times$ canal et type $\times$ canal existent ; les volumes par couple sont du même ordre de grandeur.
\item \textbf{Analyse temporelle.} Volume mensuel aligné sur le nombre de jours du mois ; rythme proche de 548 campagnes/jour.
\end{itemize}

\begin{lstlisting}[language=SQL, caption={Exemple d'agrégation SQL (ROI moyen par type)}, label={lst:sql-groupby}]
SELECT
    campaign_type,
    COUNT(*) AS campaign_count,
    AVG(roi) AS avg_roi
FROM public.marketing_campaigns
GROUP BY campaign_type
ORDER BY avg_roi DESC;
\end{lstlisting}

Résultat observé : Influencer 40\,170 campagnes (ROI moyen 5,0110), Search 40\,157 (5,0084), Display 39\,988 (5,0065), Email 39\,871 (4,9943), Social Media 39\,819 (4,9918). La somme des effectifs vaut 200\,005. Les écarts de ROI moyen restent inférieurs à 0,02 : ce n'est pas un classement de performance.

Des tableaux SQL complémentaires figurent en annexe~\ref{app:sql}.

\chapter{Extraction et préparation Python}
\label{chap:python}

\section{Connexion et extraction}

Le script \filename{postgres\_to\_python.py} établit une connexion SQLAlchemy \texttt{postgresql+psycopg2} vers \texttt{127.0.0.1:5432/marketing\_analysis}, avec le rôle \texttt{data\_project}. Le mot de passe est demandé par \texttt{getpass} et n'est jamais inscrit dans le code \cite{sqlalchemydocs}. La requête est :

\begin{lstlisting}[language=SQL, caption={Extraction complète, lecture seule}, label={lst:select-star}]
SELECT *
FROM public.marketing_campaigns
ORDER BY campaign_id;
\end{lstlisting}

Le résultat est chargé par \texttt{pandas.read\_sql}. Le moteur est fermé dans un \texttt{finally} (\texttt{engine.dispose()}). Shape à l'extraction : $(200005,\ 15)$.

\section{Contrôle des types}

À l'arrivée dans Pandas : entiers en \texttt{int64}, numériques en \texttt{float64}, catégories en \texttt{str}, \texttt{campaign\_date} d'abord en \texttt{object}. \texttt{df.info()} indique 200\,005 non-null sur chaque colonne. Aucune colonne entièrement vide.

\section{Préparation}

Le script \filename{python\_dataframe\_profiling.py} convertit la date et ajoute six colonnes, sans modifier les 15 colonnes sources :

\begin{itemize}
\item \texttt{year}, \texttt{month}, \texttt{month\_name}, \texttt{quarter} dérivés de \texttt{campaign\_date} ;
\item \texttt{ctr} = clics / impressions, avec valeur manquante si impressions nulles (cas non observé) ;
\item \texttt{ctr\_percent} = $\texttt{ctr} \times 100$.
\end{itemize}

Le DataFrame préparé a pour shape $(200005,\ 21)$.

\section{Contrôles qualité}

\begin{table}[H]
\centering
\caption{Contrôles de préparation (Pandas)}
\label{tab:prep-checks}
\begin{tabular}{|p{7.2cm}|p{5.5cm}|}
\hline
\tableheadercell{Contrôle} & \tableheadercell{Résultat} \\
\hline
Shape source $(200005,\ 15)$ & OK \\
\hline
\alternaterow
\texttt{campaign\_id} unique, min 1, max 200\,005 & OK \\
\hline
Valeurs manquantes & 0 \\
\hline
\alternaterow
Dates dans $[2021$-$01$-$01,\ 2021$-$12$-$31]$ & OK \\
\hline
Cohérence \texttt{ctr} = clics / impressions & OK \\
\hline
\alternaterow
Colonnes sources inchangées après préparation & OK \\
\hline
\end{tabular}
\end{table}

Ces contrôles sont repris à l'issue de l'EDA (\filename{python\_eda.py}) : 0 NaN, 0 infini, 0 anomalie.

\chapter{Analyse exploratoire}
\label{chap:eda}

L'EDA est produite par \filename{python\_eda.py}, après le même chargement et la même préparation que le profilage. Les analyses A à S sont faites dans Pandas. Aucun graphique n'est créé à cette étape. Les contrôles post-EDA sont tous OK.

\section{KPI globaux}

\begin{table}[H]
\centering
\caption{KPI globaux sur 200\,005 campagnes}
\label{tab:kpi-globaux}
\scriptsize
\begin{tabular}{|p{3.3cm}|r|r|r|r|r|}
\hline
\tableheadercell{Variable} & \tableheadercell{Moyenne} & \tableheadercell{Médiane} & \tableheadercell{Min} & \tableheadercell{Max} & \tableheadercell{Écart-type} \\
\hline
ROI & 5,0024 & 5,0100 & 2,00 & 8,00 & 1,7345 \\
\hline
\alternaterow
Coût d'acquisition & 12\,504,44 & 12\,497 & 5\,000 & 20\,000 & 4\,337,66 \\
\hline
Taux de conversion & 0,0801 & 0,0800 & 0,01 & 0,15 & 0,0406 \\
\hline
\alternaterow
CTR (\%) & 14,04 & 9,98 & 1,01 & 99,20 & 13,09 \\
\hline
Clics & 549,77 & 550 & 100 & 1\,000 & --- \\
\hline
\alternaterow
Impressions & 5\,507,31 & 5\,518 & 1\,000 & 10\,000 & --- \\
\hline
Engagement & 5,49 & 5 & 1 & 10 & --- \\
\hline
\alternaterow
Durée (jours) & 37,50 & 30 & 15 & 60 & --- \\
\hline
\end{tabular}
\end{table}

Le total des coûts d'acquisition vaut 2\,500\,950\,881 (unités de la source). Le CTR moyen (14,04~\%) est nettement supérieur à la médiane (9,98~\%) : la distribution est asymétrique à droite.

\section{Entreprises, types, canaux, audiences, segments, localisations}

Pour chaque dimension, la somme des volumes vaut 200\,005. Les ROI moyens restent autour de 5. Aucun classement « meilleure / pire » n'est établi.

\begin{table}[H]
\centering
\caption{Étendue des volumes et du ROI moyen par dimension}
\label{tab:eda-categories}
\begin{tabular}{|p{3.4cm}|r|r|r|}
\hline
\tableheadercell{Dimension} & \tableheadercell{$n$ (min--max)} & \tableheadercell{ROI min--max} & \tableheadercell{Écart ROI} \\
\hline
Entreprise & 39\,711--40\,238 & 4,9914--5,0071 & 0,016 \\
\hline
\alternaterow
Type & 39\,819--40\,170 & 4,9918--5,0110 & 0,019 \\
\hline
Canal & 32\,820--33\,599 & 4,9887--5,0187 & 0,030 \\
\hline
\alternaterow
Audience & 39\,688--40\,259 & 4,9828--5,0206 & 0,038 \\
\hline
Segment & 39\,742--40\,210 & 4,9994--5,0043 & 0,005 \\
\hline
\alternaterow
Localisation & 39\,751--40\,269 & 4,9802--5,0123 & 0,032 \\
\hline
\end{tabular}
\end{table}

Conversion, CTR moyen et engagement restent du même ordre sur toutes ces dimensions (conversion $\approx$ 0,080, CTR $\approx$ 14~\%, engagement $\approx$ 5,49, coût $\approx$ 12\,500).

\section{Temps et durée}

Le rythme quotidien est d'environ \metric{548} campagnes par jour civil tous les mois (décembre : 547,35). Le volume plus bas de février (15\,344) correspond à 28 jours $\times$ 548, et non à une baisse d'activité distincte.

\begin{table}[H]
\centering
\caption{KPI par trimestre}
\label{tab:eda-trimestres}
\begin{tabular}{|l|r|r|r|r|r|}
\hline
\tableheadercell{Trimestre} & \tableheadercell{$n$} & \tableheadercell{ROI} & \tableheadercell{Coût} & \tableheadercell{Conversion} & \tableheadercell{CTR (\%)} \\
\hline
Q1 & 49\,320 & 5,0140 & 12\,489,67 & 0,0800 & 14,08 \\
\hline
\alternaterow
Q2 & 49\,868 & 5,0009 & 12\,471,62 & 0,0801 & 13,99 \\
\hline
Q3 & 50\,419 & 4,9921 & 12\,525,93 & 0,0799 & 14,00 \\
\hline
\alternaterow
Q4 & 50\,398 & 5,0029 & 12\,529,87 & 0,0802 & 14,09 \\
\hline
\end{tabular}
\end{table}

ROI mensuel : minimum 4,9764 (août), maximum 5,0222 (janvier), écart 0,046. Cette variation n'est pas lue comme une tendance.

Les quatre durées (15, 30, 45, 60 jours) représentent chacune environ 25~\% des campagnes. L'écart de ROI moyen entre durées est 0,012.

\section{Classes de ROI, CTR, conversion et engagement}

Les classes de ROI (faible $<$ 4, moyen $[4,6[$, élevé $\geq$ 6) forment trois tiers presque égaux (33,25~\%, 33,22~\%, 33,53~\%). Coût, conversion, CTR et engagement ne séparent pas ces classes.

Le CTR a pour quartiles P25 = 5,86~\% et P75 = 16,97~\%. Les classes de CTR donnent des ROI moyens de 5,0154 (CTR $<$ 5~\%), 4,9985 (5--15~\%) et 5,0004 ($\geq$ 15~\%).

Le taux de conversion prend 15 valeurs distinctes (0,01 à 0,15). Le ROI moyen par valeur reste dans 4,96--5,04, sans gradient.

Les scores d'engagement de 1 à 10 pèsent chacun environ 10~\% (19\,882 à 20\,141 campagnes). Le ROI moyen par score reste dans 4,98--5,02.

\section{Corrélations et croisements}

La grille descriptive de $|r|$ de Pearson est : $<$ 0,10 négligeable ; 0,10--0,30 faible ; 0,30--0,50 modérée ; 0,50--0,70 notable ; 0,70--0,90 forte ; $\geq$ 0,90 très forte. Une corrélation n'est pas une causalité.

Seules deux paires ont $|r| \geq 0,10$ : impressions vs CTR ($r = -0,6575$, notable) et clics vs CTR ($r = +0,5065$, notable). Ces liens sont mécaniques : $\mathrm{CTR} = \mathrm{clics}/\mathrm{impressions}$. Le ROI n'est associé linéairement ni au CTR, ni à la conversion, ni à l'engagement ($|r| < 0,10$).

Les 30 combinaisons entreprise $\times$ canal, 30 type $\times$ canal et 25 audience $\times$ segment sont toutes observées. Le ROI des cellules reste dans une bande étroite ($\approx$ 4,96--5,04).

Aucune modalité n'est inférée comme dominante ($>$ 40~\%) ou minoritaire ($<$ 5~\%). L'écart max--min de parts est au plus 0,39 point (canaux).

\chapter{Visualisations}
\label{chap:visualisations}

Les seize figures sont produites par \filename{python\_visualisations.py} (Matplotlib, sans Seaborn ni Plotly) et enregistrées dans \filename{visualisations/}. La revue n'a corrigé que \filename{04\_volume\_mensuel.png} (échelle honnête 540--556 campagnes/jour). Les graphiques ci-dessous sont ceux classés essentiels ; les autres sont en annexe~\ref{app:viz}.

\section{Volume par entreprise}

La \figref{volume-company} montre des effectifs très proches entre les cinq entreprises (environ 20~\% chacune). Les libellés indiquent l'effectif et la part. Ce graphique décrit un volume, pas une performance.

\standardfigure{assets/images/charts/01_volume_company.png}{0.92\textwidth}{Répartition des campagnes par entreprise (volumes, pas un classement de performance)}{fig:volume-company}

\section{Volume mensuel}

La \figref{volume-mensuel} superpose le volume mensuel brut et le rythme quotidien (n / nombre de jours du mois). Après correction de l'échelle, le rythme apparaît plat, autour de 548 campagnes/jour. Février n'est plus bas que parce qu'il a 28 jours.

\standardfigure{assets/images/charts/04_volume_mensuel.png}{0.92\textwidth}{Volume mensuel et campagnes par jour civil en 2021 (échelle quotidienne 540--556)}{fig:volume-mensuel}

\section{Distribution du CTR}

La \figref{dist-ctr} présente le CTR en pourcentage. La moyenne (14,04~\%) est supérieure à la médiane (9,98~\%) : asymétrie à droite. L'axe s'étend de 0 à 100~\%.

\standardfigure{assets/images/charts/07_distribution_ctr.png}{0.92\textwidth}{Distribution du CTR (\%)}{fig:dist-ctr}

\section{Corrélations}

La \figref{corr} affiche les coefficients de Pearson entre ROI, coût d'acquisition, conversion, clics, impressions, engagement, durée et CTR. Les deux coefficients visuellement saillants concernent CTR--clics et CTR--impressions.

\standardfigure{assets/images/charts/10_heatmap_correlations.png}{0.88\textwidth}{Corrélations de Pearson entre variables numériques}{fig:corr}

\begin{warningbox}[Lecture du CTR]
CTR = Clicks / Impressions. Les corrélations CTR--Clicks et CTR--Impressions sont liées à la définition du ratio et ne doivent pas être interprétées comme des relations causales.
\end{warningbox}

\section{ROI selon la classe de CTR}

La \figref{roi-ctr} compare les distributions de ROI pour CTR faible ($<$ 5~\%), moyen ($[5,15[$) et élevé ($\geq$ 15~\%). Les boîtes restent très proches ; les ROI moyens valent 5,0154, 4,9985 et 5,0004.

\standardfigure{assets/images/charts/12_roi_classes_ctr.png}{0.82\textwidth}{ROI selon la classe de CTR}{fig:roi-ctr}

\chapter{Dashboard Streamlit}
\label{chap:dashboard}

\section{Architecture}

Le livrable de restitution est \filename{dashboard.py}, lancé par \texttt{streamlit run dashboard.py}. Il s'appuie sur Streamlit, Pandas et Matplotlib uniquement \cite{streamlitdocs}. Le DataFrame source est chargé une fois depuis PostgreSQL (\texttt{@st.cache\_data}), en lecture seule. Les filtres produisent un DataFrame distinct ; le DataFrame chargé n'est jamais modifié.

\section{Filtres et KPI}

La barre latérale propose les filtres Company, Campaign Type, Channel, Target Audience, Customer Segment, Location, Quarter et Month, chacun avec l'option « Tous ». Les filtres se combinent. À chaque changement, les KPI et les graphiques concernés sont recalculés sur le sous-ensemble.

Les cartes d'indicateurs affichent : nombre de campagnes, ROI moyen, coût d'acquisition moyen, taux de conversion moyen (en pourcentage), CTR moyen (en pourcentage, par exemple 14,04~\% et non 0,1404) et engagement moyen.

\section{Organisation des vues}

Le dashboard reprend la logique des visualisations validées, recalculée dynamiquement :

\begin{enumerate}
\item vue d'ensemble (volumes entreprise, type, canal) ;
\item distributions (ROI, coût, conversion, CTR) ;
\item audience, segments et canaux (heatmaps ROI et classes de CTR) ;
\item analyse temporelle (volume mensuel et ROI mensuel, avec échelle prudente) ;
\item corrélations et analyses complémentaires (durée, engagement) ;
\item constats et limites issus de l'EDA.
\end{enumerate}

Un export CSV optionnel, nommé \filename{marketing\_campaigns\_filtered.csv}, est généré uniquement à partir du DataFrame filtré.

\section{Emplacement pour une capture}

Aucune capture d'écran du dashboard n'est jointe à ce rapport. L'encadré suivant réserve un emplacement pour une future illustration.

\begin{figure}[H]
\centering
\fbox{%
\begin{minipage}{0.88\textwidth}
\centering
\vspace{1.6cm}
\textit{Emplacement réservé --- capture d'écran du dashboard Streamlit}\\[0.3cm]
{\small Commande : \texttt{streamlit run dashboard.py}}
\vspace{1.6cm}
\end{minipage}}
\caption{Emplacement prévu pour une capture du dashboard (non disponible au moment de la rédaction)}
\label{fig:dashboard-placeholder}
\end{figure}

\chapter{Résultats et constats}
\label{chap:resultats}

Les constats ci-dessous sont uniquement ceux établis par l'EDA et confirmés par les visualisations revues. Ils décrivent le jeu étudié.

\begin{keyfindings}
\begin{itemize}
\item Les volumes sont globalement équilibrés entre entreprises, types de campagnes, audiences, localisations et segments (écarts de parts inférieurs à 0,4 point).
\item Le rythme est d'environ 548 campagnes par jour civil.
\item Février est plus faible en volume brut en raison de ses 28 jours, pas d'une baisse d'activité journalière distincte.
\item Le CTR moyen vaut 14,04~\% contre une médiane de 9,98~\% (asymétrie à droite).
\item Le ROI moyen est $\approx$ 5 (5,0024 ; médiane 5,0100) dans les différentes dimensions analysées.
\item Les écarts de ROI moyen entre catégories restent faibles (0,005 à 0,038 selon la dimension).
\item Les corrélations CTR/clicks ($r = +0,51$) et CTR/impressions ($r = -0,66$) sont liées à la définition du CTR.
\item Il n'apparaît pas de gradient clair du ROI selon le CTR, la durée ou l'engagement.
\end{itemize}
\end{keyfindings}

\begin{warningbox}[Pas de causalité]
Ces résultats sont descriptifs. Une corrélation, un écart de moyennes ou une différence de volume n'établit pas de mécanisme causal.
\end{warningbox}

\chapter{Limites de l'analyse}
\label{chap:limites}

\begin{constraints}
\begin{itemize}
\item Les données sont observationnelles.
\item L'analyse décrit des associations et des distributions ; elle ne permet pas d'établir une causalité.
\item Le jeu est très équilibré : les volumes par modalité sont proches, ce qui limite la visibilité d'éventuels contrastes.
\item Les différences de KPI entre catégories sont souvent très faibles ; une échelle trop zoomée les exagérerait.
\item Les corrélations CTR / clics / impressions sont en partie mécaniques, car le CTR est un ratio de ces deux grandeurs.
\item Les résultats sont limités au dataset étudié (200\,005 campagnes de 2021) et ne doivent pas être généralisés au-delà.
\item Le CTR moyen des campagnes (moyenne des ratios) n'est pas égal au CTR global ($\sum$ clics / $\sum$ impressions) ; les deux indicateurs répondent à des questions différentes.
\end{itemize}
\end{constraints}

Aucune recommandation d'allocation budgétaire, de ciblage ou de choix de canal n'est formulée à partir de ces écarts.


\chapter{Conclusion}
\label{chap:conclusion}
Le projet a mis en place une chaîne complète, reproductible, allant du fichier Excel à un dashboard Python, en passant par PostgreSQL, SQL, Pandas, une EDA et seize visualisations Matplotlib revues.

La chaîne technique est volontairement simple : une table relationnelle unique, des lectures SQL, une préparation Pandas bornée (calendrier et CTR), des analyses descriptives, des graphiques honnêtes sur les échelles, une restitution Streamlit filtrable. PostgreSQL et les données sources n'ont pas été altérés après l'import.

Les enseignements portent sur la structure du jeu plus que sur des leviers d'action : volumes réguliers, ROI moyen proche de 5 quelle que soit la dimension, CTR asymétrique, corrélations du CTR largement définitionnelles, absence de gradient clair du ROI selon le CTR, la durée ou l'engagement.

L'intérêt du pipeline est de séparer clairement stockage, exploration, préparation, analyse, visualisation et restitution, et de conserver une traçabilité des contrôles de qualité (0 anomalie détectée).

Des évolutions possibles, non réalisées ici, consisteraient à ajouter une capture d'écran du dashboard dans ce rapport, à documenter d'autres jeux de données s'ils devenaient disponibles, ou à étendre la restitution sans changer le caractère descriptif de l'analyse. Aucune recommandation métier n'est déduite des écarts observés.


\appendix
\chapter{Contrôles qualité}
\label{app:qualite}

\begin{table}[H]
\centering
\caption{Synthèse des contrôles de qualité}
\label{tab:annexe-qualite}
\begin{tabular}{|p{7.4cm}|p{5.4cm}|}
\hline
\tableheadercell{Contrôle} & \tableheadercell{Résultat} \\
\hline
Nombre de lignes & 200\,005 \\
\hline
\alternaterow
Colonnes sources / après préparation & 15 / 21 \\
\hline
\texttt{campaign\_id} unique & OK (1 à 200\,005) \\
\hline
\alternaterow
Valeurs manquantes & 0 \\
\hline
Infinis & 0 \\
\hline
\alternaterow
Dates hors 2021 & 0 \\
\hline
CTR cohérent avec clics / impressions & OK \\
\hline
\alternaterow
Valeurs sources inchangées après EDA & OK \\
\hline
Anomalies de qualité & 0 \\
\hline
\end{tabular}
\end{table}

\chapter{Tableaux SQL complémentaires}
\label{app:sql}

\section{Volumes par entreprise}

\begin{table}[H]
\centering
\caption{Campagnes par entreprise (SQL / briefing)}
\label{tab:annexe-company}
\begin{tabular}{|l|r|r|r|}
\hline
\tableheadercell{Entreprise} & \tableheadercell{$n$} & \tableheadercell{Part} & \tableheadercell{ROI moyen} \\
\hline
TechCorp & 40\,238 & 20,12~\% & 5,01 \\
\hline
\alternaterow
Alpha Innovations & 40\,051 & 20,02~\% & 5,01 \\
\hline
DataTech Solutions & 40\,014 & 20,01~\% & 5,01 \\
\hline
\alternaterow
NexGen Systems & 39\,991 & 20,00~\% & 4,99 \\
\hline
Innovate Industries & 39\,711 & 19,86~\% & 5,00 \\
\hline
\end{tabular}
\end{table}

\section{Classes de ROI (\texttt{CASE})}

\begin{table}[H]
\centering
\caption{Classes de ROI calculées en SQL}
\label{tab:annexe-roi-case}
\begin{tabular}{|l|r|r|}
\hline
\tableheadercell{Classe} & \tableheadercell{Campagnes} & \tableheadercell{ROI moyen} \\
\hline
Faible (ROI $<$ 4) & 66\,503 & 2,9946 \\
\hline
\alternaterow
Moyen ($4 \leq$ ROI $<$ 6) & 66\,438 & 4,9983 \\
\hline
Élevé (ROI $\geq$ 6) & 67\,064 & 6,9975 \\
\hline
\end{tabular}
\end{table}

\section{Types de campagnes}

\begin{table}[H]
\centering
\caption{Effectif et ROI moyen par type de campagne}
\label{tab:annexe-types}
\begin{tabular}{|l|r|r|}
\hline
\tableheadercell{Type} & \tableheadercell{Campagnes} & \tableheadercell{ROI moyen} \\
\hline
Influencer & 40\,170 & 5,0110 \\
\hline
\alternaterow
Search & 40\,157 & 5,0084 \\
\hline
Display & 39\,988 & 5,0065 \\
\hline
\alternaterow
Email & 39\,871 & 4,9943 \\
\hline
Social Media & 39\,819 & 4,9918 \\
\hline
\end{tabular}
\end{table}

\chapter{Visualisations complémentaires}
\label{app:viz}

Les figures suivantes complètent le chapitre~\ref{chap:visualisations}. Elles proviennent du même script Matplotlib et de la même revue.

\standardfigure{assets/images/charts/02_volume_campaign_type.png}{0.88\textwidth}{Volume par type de campagne}{fig:vol-type}

\standardfigure{assets/images/charts/03_volume_channel.png}{0.88\textwidth}{Volume par canal}{fig:vol-channel}

\standardfigure{assets/images/charts/05_roi_mensuel.png}{0.88\textwidth}{ROI moyen mensuel (écarts faibles)}{fig:roi-mois}

\standardfigure{assets/images/charts/06_distribution_roi.png}{0.88\textwidth}{Distribution du ROI}{fig:dist-roi}

\standardfigure{assets/images/charts/08_distribution_acquisition_cost.png}{0.88\textwidth}{Distribution du coût d'acquisition}{fig:dist-cout}

\standardfigure{assets/images/charts/09_distribution_conversion_rate.png}{0.88\textwidth}{Distribution discrète du taux de conversion}{fig:dist-conv}

\standardfigure{assets/images/charts/11_roi_classes.png}{0.80\textwidth}{Distribution du ROI par classe de ROI}{fig:roi-classes}

\standardfigure{assets/images/charts/13_heatmap_audience_segment.png}{0.92\textwidth}{ROI moyen : audience $\times$ segment}{fig:heat-aud}

\begin{center}
\textit{Les écarts observés sont faibles ; ils ne permettent pas d'établir une différence forte entre les combinaisons.}
\end{center}

\standardfigure{assets/images/charts/14_heatmap_company_channel.png}{0.92\textwidth}{ROI moyen : entreprise $\times$ canal (pas un classement)}{fig:heat-co}

\standardfigure{assets/images/charts/15_roi_duration.png}{0.80\textwidth}{ROI selon la durée de campagne}{fig:roi-duree}

\standardfigure{assets/images/charts/16_roi_engagement.png}{0.88\textwidth}{ROI moyen selon le score d'engagement}{fig:roi-eng}

\chapter{Structure du projet et scripts}
\label{app:scripts}

\begin{table}[H]
\centering
\caption{Scripts analytiques du projet}
\label{tab:scripts}
\begin{tabular}{|p{5.4cm}|p{8.0cm}|}
\hline
\tableheadercell{Fichier} & \tableheadercell{Rôle} \\
\hline
\texttt{postgres\_to\_python.py} & Extraction PostgreSQL $\rightarrow$ Pandas \\
\hline
\alternaterow
\texttt{python\_dataframe\_profiling.py} & Profilage et préparation (21 colonnes) \\
\hline
\texttt{python\_eda.py} & EDA A--S, sans graphique \\
\hline
\alternaterow
\texttt{python\_visualisations.py} & 16 figures Matplotlib \\
\hline
\texttt{dashboard.py} & Dashboard Streamlit \\
\hline
\end{tabular}
\end{table}

Le présent rapport LaTeX se trouve dans \filename{report/}. Le modèle de mise en page provient du dépôt GitHub \texttt{sikatikenmogne/latex-report-template}, licence MIT \cite{sikatitemplate}. Le dépôt amont n'a pas été modifié : seule une copie de travail a été adaptée.

Compilation : \texttt{python scripts/compile.py} depuis \filename{report/}, puis copie du PDF vers \filename{rapport\_analyse\_marketing.pdf}.


\chapter*{Bibliographie}
\addcontentsline{toc}{chapter}{Bibliographie}
\printbibliography[heading=none]

\end{document}
